

\section{Bucketed Vector Assembly}
\label{sec:bva}

\subsection{Aligning Witness Vectors with TBVs}
\label{sec:align}

BVA specifies a deterministic decreasing-size packing rule for arbitrary input
lengths: it defines the rounding, the final power-of-two slack, tie breaking,
and the alignment invariant used by both TBV evaluation and bucket-local
folding.

Given $T$ vectors $\{\bm{a}_j\}_{j\in[T]}$ with lengths
$\{n_j\}_{j\in[T]}$ and total size $N := \sum_{j=1}^T n_j$,
our goal is to embed them into the buckets of a TBV such that:
\begin{itemize}
    \item each vector $\bm{a}_j$ occupies one bucket; and
    \item the resulting TBV dimension $K$ remains within a constant factor of $N$.
\end{itemize}
To this end, we seek bucket sizes
$\{k_i\}_{i\in[T]}$ satisfying Definition~\ref{def:tbv},
with each $n_j$ assigned to some $k_i\ge n_j$.

For each vector, choose the power-of-two cover size
$q_j:=\mathsf{next\_power\_of\_two}(n_j)$, equivalently
$n_j\le q_j<2n_j$.  This layout rule consumes only the cover sizes
$\{q_j\}$; the exact $n_j$ values can remain committed in the size-private mode
of Section~\ref{sec:size-private}.
The vector $\bm{a}_j$ is conceptually padded with $q_j-n_j$ zeros to form one
size-$q_j$ bucket.  Let $S:=\sum_{j=1}^T q_j$ and choose
$K:=\mathsf{next\_power\_of\_two}(S)$.
The $T$ vector buckets are sorted by non-increasing cover size, with ties broken
by the original vector index.  The remaining $K-S$ coordinates form one
implicit zero suffix after all vector buckets; this suffix is not decomposed
into buckets.

\begin{proposition}[BVA packing guarantee]
\label{prop:bvabound}
For any nonempty collection of total length $N$, the decreasing-size layout rule
outputs an aligned power-of-two domain of length $K<4N$ with
exactly $T$ buckets.  Its physical representation contains
exactly the $N$ input entries plus bucket metadata; it contains no padded field
element.
\end{proposition}

\begin{proof}
The non-increasing order satisfies requirement~\textcircled{2} and guarantees
alignment without inserting padding between vector buckets: every preceding
bucket is a power of two at least as large as $k_i$, hence
$k_i\mid K_{i-1}$.  Since $q_j<2n_j$, we have $S<2N$;
rounding $S$ to the next power of two gives $K<2S<4N$.
There is exactly one bucket per input vector, so the bucket count is $T$.  Each bucket stores
only its $n_j$ input entries and length.  Neither its internal zero suffix nor
the final $K-S$ zero suffix is allocated.
\end{proof}
Figure~\ref{fig:layout} illustrates this padding and alignment procedure.

\mysmallskip
\noindent
\textbf{Zero materialized-padding overhead.}
Both sources of padding are purely conceptual: the $q_j-n_j$ zero suffix within
each vector bucket and the single $K-S$ zero suffix of the logical domain.
They are determined by the committed effective lengths $\{n_j\}$, public bucket
sizes $\{k_i\}$, and the mapping $\phi$ from input vectors to buckets.  The
public part is produced by the decreasing-size layout rule above, and
Eq.~\eqref{eq:size-private} binds the private boundaries to it.
Within a vector bucket, $n_j$ marks the end of the stored entries, so its zero
suffix is never allocated.  The final suffix is determined by $S$ and $K$ and
has no bucket record.  Subsequent operations consume vector entries and layout
metadata directly.  Thus, padding changes the logical PCS domain but adds no
stored field elements.

\begin{definition}[Bucketed Vector Assembly]
    Given $T$ vectors $\{\bm{a}_j\}_{j\in[T]}$ with lengths $\{n_j\}_{j\in[T]}$
and power-of-two covers $\bm q=(q_1,\ldots,q_T)$,
their \emph{bucketed vector assembly (BVA)} is specified by:
\[
    \mathsf{BVA}_{\bm q}(\{\bm{a}_j\}_{j\in[T]})
    :=
    \big(
        K,\,
        S,\,
        T,\,
        \{(n_j,\bm{a}_j)\}_{j\in[T]},\,
        \{k_i\}_{i\in[T]},\,
        \phi
    \big),
\]
where $K$, $S$, $\{k_i\}_{i\in[T]}$, and $\phi$ are the outputs of
the decreasing-size layout rule above.  We omit the subscript $\bm q$ when the covers are
clear from context.
The set of all such BVAs over $\Field$ is denoted by $\BVAspace$.
\end{definition}

\begin{figure}[!t]
\noindent
\resizebox{\linewidth}{!}{
% Per-vector packing, decreasing-size BVA layout, and one EO-folding round.
\definecolor{rainbowyellow}{HTML}{FFF3CE}
\definecolor{rainbowgreen}{HTML}{D6E9D5}
\definecolor{rainbowred}{HTML}{FAD8D4}
\definecolor{rainbowcyan}{HTML}{AFE3E6}
\definecolor{rainbowblue}{HTML}{D9E8FC}
\definecolor{rainbowgray}{HTML}{BAC8D3}
\definecolor{rainbowpurple}{HTML}{D9D2E9}
\definecolor{rainboworange}{HTML}{FCE5CD}
\definecolor{rainbowpink}{HTML}{EAD1DC}

\begin{tikzpicture}[
  x=0.30cm,
  y=0.70cm,
  >=Latex,
  font=\sffamily\scriptsize,
  bucket/.style={draw=black!75, line width=0.55pt, minimum height=0.58cm,
                 inner sep=0pt, align=center},
  padboundary/.style={draw=black!48, densely dashed, line width=0.5pt},
  groupboundary/.style={draw=black!60, line width=0.55pt},
  implicit/.style={draw=black!55, dashed, line width=0.55pt, fill=none},
  flow/.style={->, draw=black!55, line width=0.65pt}
]

% Draw one BVA bucket.  Only the first ``data'' coordinates are colored;
% the conceptual zero suffix remains blank.
\newcommand{\BVADataBucket}[5]{%
  % start, cover size, data length, fill color, vertical coordinate
  \path[fill=#4,draw=none] (#1,#5) rectangle ++(#3,0.82);
  \path[bucket,fill=none] (#1,#5) rectangle ++(#2,0.82);
}

% ---------------------------------------------------------------------------
% (a) Packed EO-LiteBullet's first reduction.  The colored input width is
% proportional to the native length.  Each four-coordinate block is compressed
% with y_0 into one full-field element; logical padding remains blank.
\node[anchor=west,font=\bfseries] at (0,7.10) {(a) Packed reduction ($\kappa=4$)};
\node[anchor=south,text=black!65] at (16,6.55)
  {native low-bit-width entries};

% Representative length 32 is compressed to length 8.
\node[anchor=east] at (-0.35,6.00) {$32$};
\path[fill=rainbowcyan,draw=none] (0,5.59) rectangle ++(32,0.82);
\path[bucket,fill=none] (0,5.59) rectangle ++(32,0.82);
\foreach \i in {1,...,31}{
  \pgfmathsetmacro{\cellx}{\i}
  \draw[black!22,line width=0.25pt] (\cellx,5.59) -- (\cellx,6.41);
}
\foreach \i in {4,8,...,28}{
  \pgfmathsetmacro{\groupx}{\i}
  \draw[groupboundary] (\groupx,5.59) -- (\groupx,6.41);
}
% Omitted vectors are arranged vertically between the first and last examples.
\node at (16,5.25) {$\vdots$};

% Representative length 4 is compressed to length 1.
% Its cover is centered with the length-32 cover above.
\node[anchor=east] at (13.65,4.50) {$4$};
\path[fill=rainbowpink,draw=none] (14,4.09) rectangle ++(4,0.82);
\path[bucket,fill=none] (14,4.09) rectangle ++(4,0.82);
\foreach \i in {1,...,3}{
  \pgfmathsetmacro{\cellx}{14+\i}
  \draw[black!22,line width=0.25pt] (\cellx,4.09) -- (\cellx,4.91);
}
% ---------------------------------------------------------------------------
% (b) N=S=25, K=32.  Only the nine real covers are buckets; the remaining
% seven coordinates form one implicit suffix.
\node[anchor=west,font=\bfseries] at (0,3.50) {(b) Decreasing-size BVA};
\node[anchor=east,font=\bfseries,align=right] at (-0.35,2.50)
  {$S=25$\\$K=32$};

% Input lengths: 8,4,4,2,2,2,1,1,1.
\BVADataBucket{0}{8}{8}{rainbowcyan}{2.09}
\BVADataBucket{8}{4}{4}{rainbowgreen}{2.09}
\BVADataBucket{12}{4}{4}{rainbowred}{2.09}
\BVADataBucket{16}{2}{2}{rainbowyellow}{2.09}
\BVADataBucket{18}{2}{2}{rainbowgray}{2.09}
\BVADataBucket{20}{2}{2}{rainbowblue}{2.09}
\BVADataBucket{22}{1}{1}{rainbowpurple}{2.09}
\BVADataBucket{23}{1}{1}{rainboworange}{2.09}
\BVADataBucket{24}{1}{1}{rainbowpink}{2.09}
\path[implicit] (25,2.09) rectangle ++(7,0.82);

% The packed vectors are not drawn separately.  The arrows point directly to
% their corresponding materialized segments in the BVA: 32 -> 8 and 4 -> 1.
\draw[flow] (16,5.55) --
  node[above,sloped,fill=white,inner sep=1pt] {$y_0$} (3.5,2.93);
\draw[flow] (16,4.05) --
  node[above,sloped,fill=white,inner sep=1pt] {$y_0$} (24.5,2.93);

\foreach \x/\s in {4/8,10/4,14/4,17/2,19/2,21/2,
                     22.5/1,23.5/1,24.5/1}
  \node[anchor=north,text=black!65] at (\x,2.03) {$\s$};
\node[anchor=north,text=black!65] at (28.5,2.03) {$K-S=7$};
\node[anchor=west,text=black!65] at (0,1.55)
  {$k_1\ge k_2\ge\cdots\ge k_9$; blank regions are conceptual zeros};

\draw[flow] (16,1.95) -- node[right,fill=white,inner sep=1pt] {EO-fold} (16,1.05);

% ---------------------------------------------------------------------------
% (c) K/2=16 and S'=13 after one EO fold.  The three unit buckets yield two
% units: the first pair folds together, and the last unit folds with zero.
\node[anchor=west,font=\bfseries] at (0,1.16) {(c) Folded BVA};
\node[anchor=east,font=\bfseries,align=right] at (-0.35,0.50)
  {$S'=13$\\$K/2=16$};
\BVADataBucket{0}{4}{4}{rainbowcyan}{0.09}
\BVADataBucket{4}{2}{2}{rainbowgreen}{0.09}
\BVADataBucket{6}{2}{2}{rainbowred}{0.09}
\BVADataBucket{8}{1}{1}{rainbowyellow}{0.09}
\BVADataBucket{9}{1}{1}{rainbowgray}{0.09}
\BVADataBucket{10}{1}{1}{rainbowblue}{0.09}
\BVADataBucket{11}{1}{1}{rainbowpurple!50!rainboworange}{0.09}
\BVADataBucket{12}{1}{1}{rainbowpink}{0.09}
\path[implicit] (13,0.09) rectangle ++(3,0.82);
\foreach \x/\s in {2/4,5/2,7/2,8.5/1,9.5/1,10.5/1,
                     11.5/1,12.5/1}
  \node[anchor=north,text=black!65] at (\x,0.03) {$\s$};
\node[anchor=north,text=black!65] at (14.5,0.03) {$3$};
\node[anchor=west,text=black!65,align=left] at (16.7,0.50)
  {pair two unit buckets\\pair the last unit with implicit zero};
\end{tikzpicture}

}
\caption{Packed reduction and decreasing-size BVA layout.  (a) With
$\kappa=4$, challenge $y_0$ compresses each consecutive four-coordinate block
of the native low-bit-width input into one field element.  The representative
lengths $32$ and $4$ reduce to materialized lengths $8$ and $1$, respectively,
matching the first and last
lengths in the BVA; the arrows point directly to these BVA segments.  (b) For input
lengths $(8,4,4,2,2,2,1,1,1)$, these nine vectors form the BVA buckets.  (c) One EO
fold preserves the layout while leaving the final zero suffix implicit.}
\Description{BVA}
\label{fig:layout}
\end{figure}

\begin{algorithm}[!t]
\caption{Linear Operations for a Decreasing-Size BVA}
\label{alg:bva}
\KwIn{
$\bm{a}=\left(K,S,T,\{(n_j,\bm{a}_j)\}_{j\in[T]},
\{k_i\}_{i\in[T]},\phi\right)\in\BVAspace$
with $K\ge2$ for splitting or folding;\\
\qquad optional $y\in\Field$ (for folding);\\
\qquad optional $\bm{\zeta}=\big(K,T,\{(k_i,\eta_i,\bm{z}_i)\}_{i\in[T]}\big)
\in\TBV(\Field)$ (for inner product)
}
\KwOut{
(\textsf{Even-Odd Split}) $(\bm{a}_O,\bm{a}_E)\in\BVAspace^2$;\\
\qquad (\textsf{Even-Odd Fold}) $\bm{a}'\in\BVAspace$;\\
\qquad (\textsf{Inner Product with TBV}) $c=\langle\bm{a},\bm{\zeta}\rangle$
}

\SetKwFunction{Bucket}{Bucket}
\SetKwFunction{Assemble}{AssembleBuckets}
\SetKwFunction{EOFold}{EO-Fold}
\SetKw{KwSplit}{Even-Odd Split}
\SetKw{KwFold}{Even-Odd Fold}
\SetKw{KwIP}{Inner Product}

\BlankLine
The function $\Assemble$ preserves the order of its real-bucket list, infers
its active cover length from the bucket sizes, and represents the remainder of
the specified logical domain by one implicit zero suffix; it stores only the
effective entries of each result bucket\;
Let $\widehat{\bm{a}}^{(i)}\leftarrow\Bucket(\bm{a},i)\in\Field^{k_i}$ denote
the unique real bucket~$i=\phi(j)$, whose conceptual content is
$\bm{a}_j\|\bm{0}^{k_i-n_j}$; its zero suffix is handled implicitly\;
Let $T_1\leftarrow|\{i\in[T]:k_i=1\}|$ and $J\leftarrow T-T_1$\;
\tcp*{the unit buckets are $J+1,\ldots,T$}

\BlankLine
\textbf{\KwSplit: produce $\bm{a}_O,\bm{a}_E$ of length $K/2$.}\\
Initialize $\mathsf{listO}\leftarrow[\,]$ and $\mathsf{listE}\leftarrow[\,]$\;
\For{$i\leftarrow1$ \KwTo $J$}{
    $\mathsf{listO}.\mathsf{append}\big(\tfrac{k_i}{2},
      \widehat{\bm{a}}^{(i)}_O\big)$\;
    $\mathsf{listE}.\mathsf{append}\big(\tfrac{k_i}{2},
      \widehat{\bm{a}}^{(i)}_E\big)$\;
}
\For{$r\leftarrow1$ \KwTo $\lfloor T_1/2\rfloor$}{
    $p\leftarrow J+2r-1$\;
    $\mathsf{listO}.\mathsf{append}(1,\widehat{\bm{a}}^{(p)}[0])$\;
    $\mathsf{listE}.\mathsf{append}(1,\widehat{\bm{a}}^{(p+1)}[0])$\;
}
\If{$T_1$ is odd}{
    $\mathsf{listO}.\mathsf{append}(1,\widehat{\bm{a}}^{(T)}[0])$
    \tcp*{its odd partner is in the implicit zero suffix}
}
$\bm{a}_O\leftarrow\Assemble(K/2,\mathsf{listO})$\;
$\bm{a}_E\leftarrow\Assemble(K/2,\mathsf{listE})$\;
\Return $(\bm{a}_O,\bm{a}_E)$\;

\BlankLine
\textbf{\KwFold: compute $\bm{a}'=\bm{a}_O+y\bm{a}_E$.}\\
Initialize $\mathsf{listFold}\leftarrow[\,]$\;
\For{$i\leftarrow1$ \KwTo $J$}{
    $\mathsf{listFold}.\mathsf{append}\big(\tfrac{k_i}{2},
      \EOFold(\widehat{\bm{a}}^{(i)},y)\big)$\;
}
\For{$r\leftarrow1$ \KwTo $\lfloor T_1/2\rfloor$}{
    $p\leftarrow J+2r-1$\;
    $\mathsf{listFold}.\mathsf{append}\big(1,
      \widehat{\bm{a}}^{(p)}[0]+y\widehat{\bm{a}}^{(p+1)}[0]\big)$\;
}
\If{$T_1$ is odd}{
    $\mathsf{listFold}.\mathsf{append}(1,\widehat{\bm{a}}^{(T)}[0])$
    \tcp*{$\widehat{\bm a}^{(T)}[0]+y\cdot0$}
}
$\bm{a}'\leftarrow\Assemble(K/2,\mathsf{listFold})$\;
\Return $\bm{a}'$\;

\BlankLine
\textbf{\KwIP: compute $c=\langle\bm{a},\bm{\zeta}\rangle$.}\\
$c\leftarrow0$\;
\For{$j\leftarrow1$ \KwTo $T$}{
    $i\leftarrow\phi(j)$\;
    $\bm{v}\leftarrow\bm{a}_j\|\bm{0}^{k_i-n_j}$
    \tcp*{conceptual padding only}
    \For{$s\leftarrow1$ \KwTo $\log_2k_i$}{
        $\bm{v}\leftarrow\EOFold(\bm{v},z_{i,s})$\;
    }
    $c\leftarrow c+\eta_i\bm{v}[0]$\;
}
\Return $c$\;
\end{algorithm}


\subsection{Bucket-Wise Packed EO-LiteBullet Proof}
\label{sec:bucketpackedeolitebullet}

BVAs are naturally compatible with EO splitting and folding.  Let
$\widehat{\bm a}^{(i)}\in\Field^{k_i}$ be the conceptual content of real
bucket~$i$, including only its intra-bucket zero suffix, and let
$T_1:=|\{i:k_i=1\}|$ and $J:=T-T_1$.  For every $i\le J$, EO folding is
performed locally:
\[
  \widehat{\bm a}^{(i)}
  \longmapsto
  \EOfold_y(\widehat{\bm a}^{(i)})
  :=\widehat{\bm a}^{(i)}_O+y\widehat{\bm a}^{(i)}_E,
\]
producing a bucket of size $k_i/2$.  The unit buckets form a contiguous suffix
of the real buckets.  Consecutive units are paired in their existing order; if
$T_1$ is odd, then $S$ is odd and, for $K\ge2$, $S<K$; the last unit is paired with the
first coordinate of the implicit zero suffix and therefore remains unchanged.
The folded cover sizes are
\[
  \bigl(k_1/2,\ldots,k_J/2,
  \underbrace{1,\ldots,1}_{\lceil T_1/2\rceil}\bigr).
\]
They are non-increasing and sum to
$S'=\lceil S/2\rceil$.  The folded logical domain has length $K/2$, with one
implicit trailing zero suffix of length $K/2-S'$.  Thus, EO folding preserves
the decreasing-size alignment invariant without inserting, sorting, or
materializing any padding bucket.  Across all rounds, each real bucket is folded
locally until it becomes a unit bucket and is then merged at most once per
round with its adjacent unit.
An illustration of EO-folding on a BVA is shown in
Figure~\ref{fig:layout}.

Finally, consider a BVA
\(
\bm{a}=\left(
    K, S, T, \{(n_j, \bm{a}_j)\}_{j\in [T]}, \{k_i\}_{i\in [T]}, \phi
\right) \in \BVAspace
\)
and a TBV
\(
\bm{\zeta}
=
\big(K,T,\{(k_i,\eta_i,\bm{z}_i)\}_{i\in[T]}\big)
\in\TBV(\Field)
\)
with matching bucket sizes.
Their inner product can be computed bucket by bucket:
\begin{equation}
\label{eq:ipbvatbv}
    \langle \bm{a}, \bm{\zeta} \rangle
    =
\sum_{j=1}^T
    \eta_{\phi(j)} \cdot \tensorproj{\bm{a}_j}{\bm{z}_{\phi(j)}},
\end{equation}
where the final zero suffix contributes nothing.

Since EO-splitting, EO-folding, and inner-product computation between a BVA and a
TBV are all local operations and do not require materializing padded buckets,
the \packedeolitebullet protocol in
Algorithm~\ref{alg:packedeolitebullet} can be executed in
a streaming, bucket-by-bucket manner.
Moreover, EO-folding within each bucket can be performed in place.
Concretely, if $\{\bm{a}_j\}_{j\in[T]}$ are stored as $N$ field elements,
EO-folding adds each even-indexed element to its adjacent odd-indexed element
and then discards the even-indexed elements.
The packing-layer auxiliary memory required for the BVA and TBV metadata is
$\Theta(T\log K)=\Theta(T\log N)$.
It contains no term proportional to the logical domain length $K$.  Therefore,
when the inputs occupy $N$ field elements, total memory is $O(N+T\log N)$.
For the binary specialization $\kappa=2$, the prover performs approximately
$3N$ field multiplications, plus $O(T)$ work for TBV folding.  This count
depends on stored entries and bucket metadata, not on the number of conceptual
zero coordinates in $K$.

\subsection{Native Short-Bit-Width Witnesses}
\label{sec:short-width}

Suppose each witness entry is a $b$-bit integer represented canonically in
$\Field$.  Storing the BVA as $N$ field words would introduce a second,
orthogonal expansion even though no zero coordinate is allocated.  The packed
round avoids this expansion.  The prover streams each contiguous
stride-$\kappa$ segment from the native $b$-bit input, embeds its entries into
$\Field$ only while accumulating the segment polynomials, and writes one folded
field element.  It never constructs an $N$-element field array.

This streaming order remains valid at bucket boundaries.  Every bucket of size
at least $\kappa$ begins at a $\kappa$-aligned offset and splits into contiguous
$\kappa$-element segments.  Smaller buckets form a suffix; the prover consumes
that suffix in global adjacent order, allowing a segment to cross bucket
boundaries while consulting the layout for conceptual zeros. Equivalently,
this is $\tau=\log_2\kappa$ applications of the BVA EO-folding rule with
correlated challenges. The lane vectors in
Section~\ref{sec:packedeolitebullet} are only an algebraic transpose of this
blockwise scan; the prover never materializes or scans them as separate
residue-class arrays.

Let $M(\kappa)$ denote the field-operation cost of multiplying two
degree-$(\kappa-1)$ polynomials.  The packed round takes
$O((N/\kappa)M(\kappa)+I)$ prover work and leaves
$O(N/\kappa+\kappa+I)$ full-field values and metadata; segments lying entirely
in the implicit final suffix are skipped.
Subsequent binary EO rounds take $O(N/\kappa+I)$ work.  For a fixed machine
representation and fixed $\kappa$, the total is linear in the number of real
entries.  Choosing a power of two $\kappa$ near the ratio between a field word
and $b$ bits makes the reduced full-field state the same asymptotic size as the
native input.

ParComm therefore materializes neither form of representation padding: it
allocates neither zero coordinates introduced by length alignment nor an
$N$-element field-width expansion of short inputs.
